\chapter*{Introduction}

% Ajouter manuellement cette section à la table des matières
\addcontentsline{toc}{chapter}{Introduction}


La remontée rapide des taux d'intérêt amorcée en 2022 par la Banque Centrale Européenne a profondément reconfiguré l'environnement dans lequel opèrent les assureurs vie français. Après une décennie de taux bas, voire négatifs, qui avait comprimé les rendements obligataires et réduit les marges des fonds en euros, le cycle de hausse a créé une nouvelle tension : les portefeuilles constitués à des taux faibles se sont retrouvés en situation de moins-values latentes importantes, tandis que les assurés, attirés par la remontée des taux sans risque sur d'autres supports d'épargne, ont exercé une pression accrue sur les rachats. Ce retournement de conjoncture a mis en lumière l'importance d'outils de modélisation prospective capables d'anticiper les effets de différents scénarios de taux sur la situation bilancielle des fonds en euros.

C'est dans ce contexte que s'inscrit le présent mémoire, réalisé au sein de la practice \textit{Insurance} d'Accenture France.

Accenture est une entreprise mondiale de services professionnels spécialisée dans le conseil en stratégie et management, la technologie et l'externalisation. Fondée en 1989 et introduite à la Bourse de New York (NYSE : ACN), elle emploie aujourd'hui plus de 700\,000 collaborateurs répartis dans plus de 120 pays. En France, Accenture est présente depuis plusieurs décennies et intervient auprès de grandes entreprises dans l'ensemble des secteurs particulièrement les secteurs financier et assurantiel.

La construction d'un générateur de scénarios économiques (GSE) monde réel constitue un socle technique pour les modèles ALM projetant les rendements réels des actifs. Ainsi, ce GSE doit simuler des trajectoires plausibles de taux d'intérêt, d'inflation, de rendements actions et immobilier sur un horizon de projection long, typiquement 40 ans, afin d'alimenter le modèle ALM et d'en produire les indicateurs de sortie. Ces indicateurs sont généralement les taux de rendement des actifs, le taux servi aux assurés, le niveau de provision pour participation aux excédents, le taux de rachat et les différentes marges. Or le paramétrage de ce GSE, en particulier le choix des niveaux de long terme et des vitesses de retour à la moyenne des processus de taux, a une influence déterminante sur les résultats du modèle ALM et sur les décisions de gestion qui en découlent.

La problématique centrale de ce mémoire est la suivante : dans quelle mesure les hypothèses retenues pour calibrer un GSE monde réel, qu'elles soient fondées sur l'estimation statistique des données historiques ou fixées à dire d'expert, conditionnent-elles les résultats du modèle ALM, et quel est l'impact de scénarios de taux choqués sur la situation d'un fonds en euros ?

Le premier chapitre pose les fondements théoriques et pratiques du travail : il présente les caractéristiques du fonds en euros, les principes du modèle ALM retenu et les modèles stochastiques utilisés dans le GSE. Le deuxième chapitre décrit le GSE monde réel mis en place : son architecture, les modèles de taux et d'actifs risqués et la génération des scénarios Monte-Carlo. Le troisième chapitre est consacré à la détermination des paramètres du GSE et à l'analyse des résultats du modèle ALM sous deux scénarios témoins, l'un calibré par moindres carrés ordinaires sur données historiques, l'autre fixé à dire d'expert selon les objectifs de la BCE. Le quatrième chapitre explore enfin la sensibilité de ces résultats à des modifications ciblées des paramètres du GSE, au travers de scénarios choqués en volatilité, en niveau de long terme et en vitesse de retour à la moyenne.
